\begin{figure*}[t]
\centering
\includegraphics[width=0.95\linewidth]{Figures_Final/adapt4.pdf}
\caption{(a) The off-road race track designed for online adaptation~(\textbf{Q3}). The vehicle trajectories are visualized, and the colors of the lines illustrate the rotational impacts exerted on the vehicle. (b) Vertical acceleration of the vehicle during navigation. By conducting online adaptation, the vehicle can plan paths that can minimize impacts exerted on the vehicle, leading to stable navigation in off-road.
}
\label{fig:adaptation}
\vspace*{-0.15in}
\end{figure*}


Our previous works have demonstrated our methods for minimizing epistemic uncertainty in vehicle dynamics and terrain traversability~\cite{kim2023bridging,seo_scate_2023}. Therefore, we concentrate experiments on the integration of methods for autonomous navigation, particularly in terms of online adaptation for traversability~(see Section~\ref{sec:meta}). Our experiments address the following key questions: \textbf{(Q1)} Can our method reduce aleatoric uncertainty for learning traversability cost? \textbf{(Q2)} Is our method helpful for safe and stable navigation in unstructured environments? \textbf{(Q3)} Can our traversability estimation model online adapt effectively to unknown terrains for navigation?
